\subsection*{示例代码：带分块的矩阵乘法}

\begin{lstlisting}[language=croqtile,caption={示例代码：带分块的矩阵乘法}]
__co__ s32 [512, 1024] matmul(s32 [512, 1024] lhs, s32 [1024, 1024] rhs) {
  // Device program
  s32[lhs.span(0), rhs.span(1)] output; // Use same shape as lhs
  parallel p by 1 {
    with index = {m_tile, k_tile, n_tile} in [4, 4, 8] {
      foreach m_tile, n_tile {
        shared s32[128, 128] l2_out;
        foreach k_tile {
          // Initiate DMA operations
          lhs_load = dma.copy.async lhs.chunkat(m_tile, k_tile) => shared;
          rhs_load = dma.copy.async rhs.chunkat(k_tile, n_tile) => shared;
          out_store_s = dma.any;
          parallel q by 1 {
            with index = {m_tile_s, k_tile_s, n_tile_s} in [4, 4, 8] {
              // Local buffers for tiling
              local s32 [32, 64] lhs_load_s_buffer_0 {0}, lhs_load_s_buffer_1 {0};
              local s32 [64, 16] rhs_load_s_buffer_0 {0}, rhs_load_s_buffer_1 {0};
              foreach index {
                // Add "after" primitive to chain up paired DMA operations
                lhs_load_s = dma.copy.async lhs_load.data.chunkat(m_tile_s, k_tile_s) => select(k_tile_s % 2, lhs_load_s_buffer_0, lhs_load_s_buffer_1) after lhs_load;
                rhs_load_s = dma.copy.async rhs_load.data.chunkat(k_tile_s, n_tile_s) => select(k_tile_s % 2, rhs_load_s_buffer_0, rhs_load_s_buffer_1) after rhs_load;
                out_load_s = dma.copy.async l2_out.chunkat(m_tile_s, n_tile_s) => local;
                wait out_load_s;
                // Perform the matrix multiplication (dot product) operation
                wait lhs_load_s, rhs_load_s;
                call dot_general_fp32(lhs_load_s.data, rhs_load_s.data, 32, 16, 64, 0, 8, out_load_s.data);
                // Store the result back to l2_out
                out_store_s = dma.copy.async out_load_s.data => l2_out.chunkat(m_tile_s, n_tile_s);
              }
            }
          }
          // Add "after" primitive to chain up DMA operations for the final output
          out_store = dma.copy.async l2_out => output.chunkat(m_tile, n_tile) after out_store_s;
          wait out_store;
        }
      }
    }
  }
  return output;
}
\end{lstlisting}

外层范围 \co@[4, 4, 8]@ 把操作数切成 \co@[128, 256]@ 的左分块、\co@[256, 128]@ 的右分块，以及 \co@[128, 128]@ 的输出分块——这正是 \co@l2_out@ 暂存缓冲区的大小。内层同样按 \co@[4, 4, 8]@ 划分，得到 \co@[32, 64]@、\co@[64, 16]@ 和 \co@[32, 16]@ 的结果。\co@after@ 子句使内层输入复制在外层复制之后开始。这段清单用于展示结构；作为可运行核函数使用前，必须补齐并检查辅助函数定义、操作数使用与缓冲区复用前的完整等待，以及输出初始化。
